% ==============================================================================
% CONFIG: PACKAGES.TEX
% Khai báo toàn bộ các gói thư viện (packages) sử dụng trong tài liệu.
% Được phân loại theo từng nhóm chức năng rõ ràng.
% ==============================================================================

% ------------------------------------------------------------------------------
% 1. ĐỊNH DẠNG KHỔ GIẤY VÀ CĂN LỀ
% ------------------------------------------------------------------------------
\usepackage[
    a4paper,
    left=3.0cm,
    right=2.0cm,
    top=2.5cm,
    bottom=2.5cm
]{geometry}

% ------------------------------------------------------------------------------
% 2. FONT CHỮ VÀ HỖ TRỢ XELATEX
% ------------------------------------------------------------------------------
\usepackage{fontspec}

% ------------------------------------------------------------------------------
% 3. CÔNG CỤ SOẠN THẢO VĂN BẢN VÀ KHOẢNG CÁCH
% ------------------------------------------------------------------------------
\usepackage{setspace}       % Điều chỉnh khoảng cách dòng
\usepackage{indentfirst}    % Thụt đầu dòng đoạn văn đầu tiên sau tiêu đề
\usepackage{enumitem}       % Tùy biến danh sách bullet, numbering
\usepackage{titlesec}       % Định dạng tiêu đề chương, mục
\usepackage{titletoc}       % Định dạng mục lục
\usepackage{fancyhdr}       % Cấu hình header và footer

% ------------------------------------------------------------------------------
% 4. TOÁN HỌC VÀ KÝ HIỆU KHOA HỌC
% ------------------------------------------------------------------------------
\usepackage{amsmath}
\usepackage{amssymb}
\usepackage{amsfonts}

% ------------------------------------------------------------------------------
% 5. HÌNH ẢNH VÀ ĐỒ HỌA
% ------------------------------------------------------------------------------
\usepackage{graphicx}       % Chèn hình ảnh
\usepackage{float}          % Quản lý vị trí hình/bảng (tùy chọn [H], [htbp])
\usepackage{subcaption}     % Chèn các hình con (subfigures)
\usepackage{tikz}           % Vẽ khung bìa và sơ đồ vector
\usetikzlibrary{calc}

% ------------------------------------------------------------------------------
% 6. BẢNG BIỂU HỌC THUẬT
% ------------------------------------------------------------------------------
\usepackage{booktabs}       % Bảng học thuật (\toprule, \midrule, \bottomrule)
\usepackage{tabularx}       % Bảng co giãn theo chiều rộng trang
\usepackage{array}          % Căn lề cột bảng nâng cao
\usepackage{longtable}      % Bảng dài qua nhiều trang
\usepackage{multirow}       % Gộp nhiều dòng trong bảng
\usepackage{adjustbox}      % Co giãn bảng vừa vặn trang giấy

% ------------------------------------------------------------------------------
% 7. ĐỊNH DẠNG CAPTION CHO HÌNH VÀ BẢNG
% ------------------------------------------------------------------------------
\usepackage{caption}

% ------------------------------------------------------------------------------
% 8. MÃ NGUỒN (SOURCE CODE LISTINGS)
% ------------------------------------------------------------------------------
\usepackage{listings}
\usepackage{xcolor}         % Quản lý màu sắc cho mã nguồn và tài liệu

% ------------------------------------------------------------------------------
% 9. TRÍCH DẪN VÀ THAM CHIẾU (IEEE STYLE)
% LaTeX mặc định đã hỗ trợ \cite{key}. BibTeX + IEEEtran sẽ định dạng số chuẩn [1].
% ------------------------------------------------------------------------------
\usepackage{url}            % Hiển thị URL chuẩn xác

% ------------------------------------------------------------------------------
% 10. LIÊN KẾT SIÊU VĂN BẢN VÀ THAM CHIẾU CHÉO (HYPERREF & CLEVEREF)
% Chú ý: hyperref nạp gần cuối, cleveref nạp sau cùng.
% ------------------------------------------------------------------------------
\usepackage[
    unicode=true,
    bookmarks=true,
    bookmarksnumbered=true,
    bookmarksopen=true,
    breaklinks=true,
    pdfstartview={FitH},
    pdfauthor={Khoa Cong nghe Thong tin - ICTU}
]{hyperref}

\usepackage[
    nameinlink,
    noabbrev,
    capitalize
]{cleveref}
